\section{Geometric bounds and determinant concentration}
\label{sec:gaud-determinant-concentration}

By \cref{thm:gaud-determinant-endpoint}, the determinant densities of all normalized Gauduchon currents are bounded in the critical space $L^{n/(n-1)}$ exactly when the lower mass is positive. Together with the mass-invariance criterion of \cref{thm:mass-invariance-characterization}, this gives geometric sufficient conditions for the critical bound. The critical exponent cannot be raised, even for currents that are powers of K\"ahler forms in one fixed class. The same concentrating family shows that currents approaching the maximum of the determinant functional for a fixed smooth density need not satisfy any uniform higher integrability, although the maximizer itself is smooth. In the subsections that follow, we compute the complete determinant profile of a background that satisfies the conditions of \cref{thm:mass-invariance-characterization}, show that this profile does not control the ellipticity of the background, and construct exact maximizers for varying densities whose weak limit has zero determinant measure.

Throughout this section $n\geq2$ and $\mathrm{d}V$ is a smooth positive probability volume form. For a Hermitian form $h$, $k_T$ denotes the determinant density relative to $\mathrm{d}V$ of a current $T\in\mathscr G_h$, as in \cref{sec:gaud-determinant-integrability}. Put $q=n/(n-1)$.

\Needspace{14\baselineskip}
\begin{corollary}\label{cor:geometric-determinant-bound}
Suppose $g$ is a Hermitian form with $\ddc g=\ddc(g^2)=0$, and put $V_g=\int_Xg^n$. For every Hermitian form $h\geq cg$, $c>0$,
\[
 v_-(h)\geq c^nV_g,\qquad
 \sup_{T\in\mathscr G_h}\|k_T\|_{L^q(\mathrm{d}V)}
       \leq c^{-1}V_g^{-1/n}.
\]
Equality holds in the determinant bound when $h=g$ and $c=1$. In particular, on a compact K\"ahler manifold or a compact complex surface, the critical norm is uniformly bounded over the normalized Gauduchon currents of each fixed Hermitian form.
\end{corollary}
\begin{proof}
The identity $\ddc(g^2)=2g\wedge\ddc g+2\mathrm{d}g\wedge\dc g$ shows that $\mathrm{d}g\wedge\dc g=0$, so \cref{thm:mass-invariance-characterization} gives $v_-(g)=V_g$. \Cref{cor:gaud-determinant-background-comparison} then gives both bounds, and \cref{thm:gaud-determinant-endpoint} gives equality in the determinant bound when $h=g$.

A K\"ahler form satisfies both differential identities. On a surface, \cref{thm:gauduchon} gives a metric $g$ with $\ddc g=0$, and $\ddc(g^2)=0$ for degree reasons. By compactness, each fixed $h$ satisfies $h\geq cg$ for some $c>0$.
\end{proof}

For a fixed smooth $F$, \cref{thm:weak-gaud-compact-duality} gives a unique maximizing current for the determinant functional $\mathfrak D_F$ of \eqref{eq:weak-gaud-functional}, and this current is smooth and strictly positive, so its determinant density is bounded. The following construction shows that currents approaching this maximum satisfy no uniform bound on higher moments of their determinant densities.

\begin{proposition}\label{prop:gaud-determinant-concentration}
Let $X=E^n$, where $E=\mathbb C/(\mathbb Z+\mathrm i\mathbb Z)$ carries its flat area form of total area one, and write $\omega_i$ for the pullback of this form from the $i$th factor. Set
\[
 h=\sum_{i=1}^n\omega_i,\qquad
 \mathrm{d}V=\omega_1\wedge\cdots\wedge\omega_n,\qquad V=n!,\qquad
 \alpha=\frac{n-1}{n}.
\]
For $F=0$, the unique maximizing current is
$T_*=h^{n-1}/V$, with $k_{T_*}=V^{-1/n}$ and
$\mathfrak D_0(T_*)=V^{-1/n}$.
There are smooth strictly positive closed forms $S_\varepsilon\in\mathscr G_h$, $0<\varepsilon<1/4$, such that
\[
 \|k_{S_\varepsilon}\|_{L^q}=V^{-1/n},\qquad
 \mathfrak D_0(S_\varepsilon)\longrightarrow V^{-1/n},
\]
while for every fixed finite $p>q$,
\begin{equation}\label{eq:gaud-concentration-moments}
 \int_X k_{S_\varepsilon}^{p}\,\mathrm{d}V
 \geq C_{n,p}[\log(1/\varepsilon)]^{-\alpha p}
                    \varepsilon^{2-2\alpha p}\longrightarrow+\infty.
\end{equation}
Their $L^\infty$ norms also diverge. Each $S_\varepsilon$ is the $(n-1)$st power of a K\"ahler form in the fixed class $V^{-1/(n-1)}[h]$, and has the same Bott--Chern, de Rham and Aeppli classes as $T_*$. No $S_\varepsilon$ attains the maximum of $\mathfrak D_0$.
\end{proposition}
\begin{proof}
Put $H_i=\bigwedge_{j\ne i}\omega_j$. Expanding the wedge powers gives
$h^{n-1}=(n-1)!\sum_iH_i$ and $h\wedge H_i=\mathrm{d}V$.
Thus $T_*=(1/n)\sum_iH_i$ is normalized, and the determinant formula gives $k_{T_*}=V^{-1/n}$. The smooth equation for $F=0$ is solved by the zero potential, since $h^n=V\,\mathrm{d}V$; \cref{thm:weak-gaud-compact-duality} therefore gives the maximal value and the uniqueness of the maximizer.

To construct $S_\varepsilon$, choose a nonnegative smooth function $\beta$ compactly supported in the unit disk in $\mathbb C$, of integral one, with $\beta\geq\beta_0>0$ on a disk of radius $r_0\in(0,1)$. On an embedded disk in the first elliptic factor put
$b_\varepsilon(z)=\varepsilon^{-2}\beta(z/\varepsilon)$, extended by zero; it has integral one. For $0<t<1$, set
\[
 \begin{gathered}
 f_{\varepsilon,t}=1-t+tb_\varepsilon,\qquad a=V^{-1/(n-1)},\\
 \gamma_{\varepsilon,t}=a(f_{\varepsilon,t}\omega_1+\omega_2+\cdots+\omega_n),
 \qquad T_{\varepsilon,t}=\gamma_{\varepsilon,t}^{n-1}.
 \end{gathered}
\]
The form $\gamma_{\varepsilon,t}$ is strictly positive, and it is closed because $f_{\varepsilon,t}$ depends only on the first factor, which has real dimension two. Direct expansion gives
\begin{equation}\label{eq:gaud-concentration-density}
 T_{\varepsilon,t}=\frac1n\left(H_1+f_{\varepsilon,t}\sum_{i=2}^nH_i\right),
 \qquad k_{T_{\varepsilon,t}}=V^{-1/n}f_{\varepsilon,t}^{\alpha}.
\end{equation}
Since $\int_Xf_{\varepsilon,t}\,\mathrm{d}V=1$, the trace integral is one, so $T_{\varepsilon,t}\in\mathscr G_h$, and
$\|k_{T_{\varepsilon,t}}\|_{L^q}=V^{-1/n}$ because $\alpha q=1$.

The mean-zero function $f_{\varepsilon,t}-1$ has a smooth real potential $\psi$ on $E$ with
$\ddc\psi=(f_{\varepsilon,t}-1)\omega_1$;
for example, the mean-zero Fourier series solving
$\Delta\psi=4\pi(f_{\varepsilon,t}-1)$ in the area-one coordinates.
Consequently $\gamma_{\varepsilon,t}-ah=\ddc(a\psi)$. Since $T_*=(ah)^{n-1}$, the closed form
$R=\sum_{j=0}^{n-2}\gamma_{\varepsilon,t}^{n-2-j}\wedge(ah)^j$ satisfies
\[
 T_{\varepsilon,t}-T_*=\ddc(a\psi R).
\]
This proves all the assertions about classes, including the case $n=2$, where $R=1$.

Jensen's inequality and $f_{\varepsilon,t}\geq1-t$ give
\[
 V^{-1/n}(1-t)^\alpha
 \leq\mathfrak D_0(T_{\varepsilon,t})
 =V^{-1/n}\int_Xf_{\varepsilon,t}^{\alpha}\,\mathrm{d}V
 <V^{-1/n},
\]
where the inequality is strict because $f_{\varepsilon,t}$ is nonconstant and $0<\alpha<1$.
On the product of the disk of radius $r_0\varepsilon$ with the other factors, we have
$f_{\varepsilon,t}\geq t\beta_0\varepsilon^{-2}$, and this set has volume $\pi r_0^2\varepsilon^2$. Hence
\[
 \begin{split}
 \int_Xk_{T_{\varepsilon,t}}^p\,\mathrm{d}V
   &\geq V^{-p/n}\pi r_0^2\beta_0^{\alpha p}t^{\alpha p}
                    \varepsilon^{2-2\alpha p},\\
 \|k_{T_{\varepsilon,t}}\|_\infty
   &\geq V^{-1/n}\beta_0^\alpha t^\alpha\varepsilon^{-2\alpha}.
 \end{split}
\]
Now take $t=1/\log(1/\varepsilon)$, which lies in $(0,1)$ because $\varepsilon<1/4<\mathrm e^{-1}$, and $S_\varepsilon=T_{\varepsilon,t}$. Then $\mathfrak D_0(S_\varepsilon)$ converges to its maximum, while both lower bounds diverge for $p>q$ because $\alpha p>1$. All constants are independent of $\varepsilon$.
\end{proof}

If instead $t\in(0,1)$ is kept fixed in the same construction, then
$\int_Xk_{T_{\varepsilon,t}}^p\,\mathrm{d}V\geq C_{n,p,t}\varepsilon^{2-2\alpha p}$ and
$\|k_{T_{\varepsilon,t}}\|_\infty\geq C_{n,t}\varepsilon^{-2\alpha}$.
Thus there are no uniform bounds on higher moments, either over the full normalized cone or along a maximizing sequence for one smooth density. Every member of these families attains the supremum $V^{-1/n}$ of the critical norm in \cref{cor:geometric-determinant-bound}, whereas the maximum of $\mathfrak D_0$ is attained only by $T_*$.
